% ----------------------------------------------------------------------------
%                                   CAPA
% ----------------------------------------------------------------------------
% Caso algum dos campos não se aplique ao seu trabalho, como por exemplo,
% se não houve coorientador, apenas deixe vazio.
% Exemplos:
% \coorientador{}
% \departamento{}

% ----------------------------------------------------------------------------
%                               DADOS DO TRABALHO
% ----------------------------------------------------------------------------
\titulo{Painel Parkinson Desktop}
\subtitulo{Migração Arquitetural para Suporte à Decisão Clínica e Otimização de Custos}
\autor{Felipe Dias Santana \\ Pedro Alexandrino Furtado Pereira Neto}
\local{Brasília}
\data{2026}
\projeto{Trabalho de Conclusão de Curso}
\tituloAcademico{Tecnólogo em Sistemas para Internet}

% ----------------------------------------------------------------------------
%                               DADOS DA INSTITUIÇÃO
% ----------------------------------------------------------------------------
\instituicao{Instituto Federal de Brasília}
\departamento{\textit{Campus} Brasília}
\programa{Tecnólogo em Sistemas para Internet}
\logoinstituicao{0.15}{dados/figuras/ifb-logo.png}

% ----------------------------------------------------------------------------
%                              DADOS DOS ORIENTADORES
% ----------------------------------------------------------------------------
\orientador[Orientador:]{Prof. Dr. Fábio Henrique Monteiro Oliveira} 
\instOrientador{Instituto Federal de Brasília}

\coorientador[Coorientador:]{Prof. Dr. João Paulo Folador}
\instCoorientador{Universidade Federal de Uberlândia}

% Quando existir mais de um coorientador
% \coorientador[Coorientadores:]{Prof. Me. XXXXXX \newline
% Universidade XXXXXX - Câmpus XXXXXX \newline
% \newline Prof. Dr. XXXXXX. 
% \newline Universidade XXXXXX - Câmpus XXXXXX}
